\documentclass[11pt,a4paper]{article}

% ---------------------------------------------------------------
% Tipografia: Palatino, versalete, uma cor de destaque (bordô)
% ---------------------------------------------------------------
\usepackage[utf8]{inputenc}
\usepackage[T1]{fontenc}
\usepackage[brazil]{babel}
\usepackage[sc,osf]{mathpazo}
\usepackage{microtype}
\usepackage[top=2.4cm, bottom=2.6cm, left=2.6cm, right=2.6cm]{geometry}
\usepackage{amsmath}
\usepackage{amssymb}
\usepackage{xcolor}
\usepackage{tcolorbox}
\tcbuselibrary{skins, breakable}
\usepackage{enumitem}
\usepackage{booktabs}
\usepackage{array}
\usepackage{tabularx}
\usepackage{needspace}
\usepackage{hyperref}

\definecolor{vinho}{RGB}{122, 27, 42}
\definecolor{grafite}{RGB}{70, 70, 70}
\hypersetup{colorlinks=true, linkcolor=vinho, urlcolor=vinho}

\linespread{1.06}
\setlength{\parindent}{0pt}
\setlength{\parskip}{0.45em}

% ---------------------------------------------------------------
% Cabeçalhos de bloco e seção
% ---------------------------------------------------------------
\newcommand{\bloco}[3]{%
  \needspace{7\baselineskip}
  \vspace{1.8em}
  {\scshape\color{vinho}bloco #1}\hfill
    {\small\itshape\color{grafite}\textasciitilde\,#3 min}\\[0.1em]
  {\Large\bfseries #2}\\[-0.55em]
  {\color{black!30}\rule{\linewidth}{0.5pt}}
  \par\nobreak\smallskip
}

\newcommand{\secao}[1]{%
  \needspace{6\baselineskip}
  \vspace{1.8em}
  {\Large\bfseries #1}\\[-0.55em]
  {\color{black!30}\rule{\linewidth}{0.5pt}}
  \par\nobreak\smallskip
}

\newcommand{\objetivo}[1]{{\itshape\color{grafite}Objetivo: #1}\par\smallskip}

% ---------------------------------------------------------------
% Painéis: barra fina à esquerda + rótulo em versalete
% ---------------------------------------------------------------
\newtcolorbox{quadro}{blanker, breakable, left=12pt,
  borderline west={1.4pt}{0pt}{grafite},
  before skip=11pt, after skip=11pt,
  before upper={{\small\scshape\color{grafite}no quadro}\par\smallskip}}

\newtcolorbox{perguntas}{blanker, breakable, left=12pt,
  borderline west={1.4pt}{0pt}{vinho},
  before skip=11pt, after skip=11pt,
  before upper={{\small\scshape\color{vinho}perguntar \`a turma}\par\smallskip}}

\newtcolorbox{obs}{blanker, breakable, left=12pt,
  borderline west={1.4pt}{0pt}{black!30},
  before skip=11pt, after skip=11pt,
  before upper={{\small\scshape\color{black!55}observa\c{c}\~oes}\par\smallskip}}

\newtcolorbox{corte}{blanker, breakable, left=12pt,
  borderline west={1.4pt}{0pt}{vinho!70, dashed},
  before skip=11pt, after skip=11pt,
  before upper={{\small\scshape\color{vinho}se o tempo apertar}\par\smallskip}}

\newtcolorbox{projetor}{blanker, breakable, left=12pt,
  borderline west={1.4pt}{0pt}{grafite},
  borderline west={1.4pt}{3.2pt}{black!25},
  before skip=11pt, after skip=11pt,
  before upper={{\small\scshape\color{grafite}projetor}\par\smallskip}}

% ---------------------------------------------------------------
\begin{document}
% ---------------------------------------------------------------

% ---------------------------------------------------------------
% TÍTULO
% ---------------------------------------------------------------
\thispagestyle{empty}

{\scshape\color{vinho}cursinho solid\'ario \,\textperiodcentered\, f\'isica}\\[-0.5em]
{\color{vinho}\rule{\linewidth}{1.6pt}}\\[0.9em]
{\huge\bfseries Din\^amica}\\[0.35em]
{\LARGE As Tr\^es Leis de Newton}\\[0.7em]
{\itshape\color{grafite}Plano de aula \,\textperiodcentered\, 2h
 \,\textperiodcentered\, ENEM e vestibular UTFPR}
\hfill {\color{grafite}2026}\\[-0.3em]
{\color{black!30}\rule{\linewidth}{0.5pt}}

\vspace{1.2em}

% ---------------------------------------------------------------
% VISÃO GERAL
% ---------------------------------------------------------------
{\renewcommand{\arraystretch}{1.45}
\begin{tabularx}{\linewidth}{@{}p{2.6cm}X@{}}
  \multicolumn{2}{@{}l}{{\scshape\color{vinho}vis\~ao geral}}\\
  \addlinespace[2pt]
  \toprule
  \textbf{Objetivo} & Enunciar as tr\^es Leis de Newton e apresentar as
    principais for\c{c}as (peso, normal, atrito e tra\c{c}\~ao),
    aplicando em quest\~oes de ENEM e UTFPR. \\
  \textbf{Formato} & Aula no quadro do in\'icio ao fim. O projetor s\'o
    entra no Bloco 5, para mostrar as quest\~oes e resolver ao vivo no
    quadro. \\
  \textbf{Materiais} & Quadro branco e canet\~oes (2 cores ou mais; uma
    s\'o para as for\c{c}as) \,\textperiodcentered\, projetor + notebook
    com o arquivo \texttt{lista\_questoes\_dinamica.pdf}
    \,\textperiodcentered\, livro, folha e caneta para as
    demonstra\c{c}\~oes \,\textperiodcentered\, lista de quest\~oes
    impressa para os alunos. \\
  \textbf{Obs.} & Esse plano junta duas aulas do planejamento original
    (Leis de Newton + For\c{c}as) em uma aula s\'o de 2h, ent\~ao o
    tempo \'e apertado. As caixas \textit{``se o tempo apertar''}
    mostram o que d\'a para cortar. \\
  \bottomrule
\end{tabularx}}

\vspace{1.4em}

% ---------------------------------------------------------------
% CRONOGRAMA
% ---------------------------------------------------------------
{\renewcommand{\arraystretch}{1.35}
\begin{tabularx}{\linewidth}{@{}clX r@{}}
  \multicolumn{4}{@{}l}{{\scshape\color{vinho}cronograma}}\\
  \addlinespace[2pt]
  \toprule
  \textbf{Bloco} & \textbf{T\'opico} & \textbf{Foco principal} & \textbf{Tempo} \\
  \midrule
  0 & Abertura & Situa\c{c}\~ao-problema: o carro que freia & 10 min \\
  1 & 1\textordfeminine{} Lei & In\'ercia, equil\'ibrio, exemplos do dia a dia & 15 min \\
  2 & 2\textordfeminine{} Lei + for\c{c}as & $F_R = ma$; peso, normal, atrito, tra\c{c}\~ao & 30 min \\
  3 & 3\textordfeminine{} Lei & A\c{c}\~ao e rea\c{c}\~ao em corpos diferentes & 15 min \\
  4 & Aplica\c{c}\~ao & Diagrama de for\c{c}as + exemplo resolvido & 15 min \\
  5 & Quest\~oes & 5 quest\~oes ENEM/UTFPR ao vivo (projetor) & 25 min \\
  6 & Encerramento & Resumo das 3 leis + tarefa & 10 min \\
  \midrule
  & \textbf{Total} & & \textbf{120 min} \\
  \bottomrule
\end{tabularx}}

% ===================================================================
\bloco{0}{Abertura e Situa\c{c}\~ao-Problema}{10}
% ===================================================================

\objetivo{come\c{c}ar com uma situa\c{c}\~ao do dia a dia (o carro que
freia e o passageiro que vai para a frente) antes de entrar na teoria.}

\begin{perguntas}
Lan\c{c}ar as perguntas oralmente e deixar os alunos responderem, sem
corrigir nada ainda:
\begin{enumerate}[nosep]
  \item ``Um carro freia bruscamente e o passageiro sem cinto vai para
        a frente. \textbf{Quem empurrou o passageiro para a frente?}''
  \item ``O que \'e preciso para um corpo ficar em movimento? Precisa
        de for\c{c}a o tempo todo?''
  \item ``Um corpo mais pesado cai mais r\'apido?''
\end{enumerate}
\textit{N\~ao dar as respostas ainda. Anotar as ideias no canto do
quadro, porque v\'arias delas aparecem como alternativas erradas nas
quest\~oes do ENEM.}
\end{perguntas}

\begin{quadro}
\begin{itemize}[nosep]
  \item Separar uma \textbf{faixa do lado direito} do quadro para ir
        montando o resumo das 3 leis durante a aula (apagar s\'o no
        final).
  \item Escrever o t\'itulo no espa\c{c}o principal:
        \textbf{Din\^amica --- por que as coisas se movem?}
  \item Anotar o que os alunos responderem sobre a pergunta do carro
        (mesmo as respostas erradas).
\end{itemize}
\end{quadro}

\begin{obs}
\textbf{Come\c{c}o da aula:} situar a turma. Na cinem\'atica a gente
s\'o descreveu o movimento, sem perguntar a causa. Agora a pergunta \'e
outra: \textbf{o que faz o movimento mudar?} A resposta s\~ao as 3 leis
de Newton, que caem muito no ENEM e na UTFPR.

\textbf{Demonstra\c{c}\~ao:} puxar uma folha de papel \textit{devagar}
com uma caneta em cima (a caneta vai junto) e depois puxar
\textit{r\'apido} (a caneta fica). \'E a mesma ideia do carro que
freia.
\end{obs}

% ===================================================================
\bloco{1}{1\textordfeminine{} Lei --- Princ\'ipio da In\'ercia}{15}
% ===================================================================

\objetivo{enunciar a 1\textordfeminine{} Lei, ligar com a
situa\c{c}\~ao do come\c{c}o e mostrar que o movimento n\~ao precisa de
for\c{c}a para continuar.}

\begin{quadro}
\textbf{1. Enunciado} (5 min)
\begin{itemize}[nosep]
  \item Escrever na faixa de resumo e ler em voz alta:
  \begin{center}
    \textbf{1\textordfeminine{} Lei:} se $\vec{F}_R = \vec{0}$, o corpo
    permanece \textbf{em repouso} ou em \textbf{MRU}.
  \end{center}
  \item Chamar aten\c{c}\~ao para o ``ou'': repouso e MRU s\~ao a mesma
        coisa do ponto de vista da din\^amica. Nenhum dos dois precisa
        de for\c{c}a para se manter.
  \item Definir \textbf{in\'ercia}: tend\^encia de manter o estado de
        movimento. Quanto mais massa, mais in\'ercia.
\end{itemize}

\medskip
\textbf{2. Voltar no carro do come\c{c}o} (5 min)
\begin{itemize}[nosep]
  \item Responder a pergunta inicial: \textbf{ningu\'em empurrou o
        passageiro}. O carro freou e o passageiro continuou em MRU por
        in\'ercia. Quem mudou de estado foi o carro.
  \item Desenhar as duas cenas: antes (carro e passageiro com
        velocidade $v$) e depois (carro parado, passageiro ainda com
        $v$).
  \item Citar outros exemplos: cinto de seguran\c{c}a, encosto de
        cabe\c{c}a, \^onibus arrancando (o passageiro ``vai para
        tr\'as'', mas na verdade ele fica).
\end{itemize}

\medskip
\textbf{3. Equil\'ibrio} (3--4 min)
\begin{itemize}[nosep]
  \item Deixar claro: $\vec{F}_R = \vec{0}$ n\~ao quer dizer que n\~ao
        tem for\c{c}a nenhuma, quer dizer que as for\c{c}as
        \textbf{se cancelam}.
  \item Exemplos no quadro: livro na mesa (peso $\downarrow$ e
        normal $\uparrow$); carro em MRU na estrada
        (motor $\rightarrow$ e atrito+ar $\leftarrow$).
\end{itemize}
\end{quadro}

\begin{perguntas}
\begin{itemize}[nosep]
  \item ``Um avi\~ao voa em linha reta com velocidade constante. A
        for\c{c}a resultante nele \'e maior, menor ou igual \`a de
        quando ele est\'a parado no p\'atio?''
        \textit{(Resposta: igual, nula nos dois casos. Muita gente
        erra.)}
  \item ``Por que na freada brusca a gente `continua' e na curva a
        gente `sai reto'?''
        \textit{(In\'ercia: o corpo tende a seguir em linha reta.)}
\end{itemize}
\end{perguntas}

\begin{obs}
Erro mais comum no ENEM: achar que in\'ercia \'e \textbf{for\c{c}a}
(``a for\c{c}a da in\'ercia jogou o passageiro''). In\'ercia \'e
propriedade do corpo, n\~ao for\c{c}a, ent\~ao n\~ao entra no diagrama
de for\c{c}as. Falar isso claramente e escrever no quadro:
\textbf{``in\'ercia n\~ao \'e for\c{c}a''}.
\end{obs}

\begin{corte}
Cortar o exemplo do avi\~ao e deixar s\'o o carro e o livro na mesa.
N\~ao cortar o enunciado nem a parte do ``ou'' (repouso/MRU).
\end{corte}

% ===================================================================
\bloco{2}{2\textordfeminine{} Lei --- $F_R = ma$ e as Principais For\c{c}as}{30}
% ===================================================================

\objetivo{enunciar a 2\textordfeminine{} Lei e apresentar as quatro
for\c{c}as que aparecem em quase toda quest\~ao de mec\^anica: peso,
normal, atrito e tra\c{c}\~ao. \'E o bloco mais pesado da aula.}

\begin{quadro}
\textbf{1. Enunciado e sentido f\'isico} (8 min)
\begin{itemize}[nosep]
  \item Escrever na faixa de resumo:
  \[
    \boxed{\vec{F}_R = m\,\vec{a}}
  \]
  \item Ler com a turma: for\c{c}a resultante produz
        \textbf{acelera\c{c}\~ao} (mudan\c{c}a de velocidade), na mesma
        dire\c{c}\~ao e sentido da for\c{c}a.
  \item Escrever as duas leituras da f\'ormula:
    \begin{itemize}[nosep]
      \item mesma for\c{c}a, massa maior $\Rightarrow$
            acelera\c{c}\~ao menor (carrinho de supermercado vazio
            vs.\ cheio);
      \item mesma massa, for\c{c}a maior $\Rightarrow$
            acelera\c{c}\~ao maior.
    \end{itemize}
  \item Unidade: $1\ \text{N} = 1\ \text{kg}\cdot\text{m/s}^2$
        (mais ou menos o peso de uma ma\c{c}\~a pequena).
  \item Ligar com a 1\textordfeminine{} Lei: se $F_R = 0$, ent\~ao
        $a = 0$. Ou seja, a 1\textordfeminine{} \'e um caso particular
        da 2\textordfeminine{}.
\end{itemize}

\medskip
\textbf{2. Peso} (5 min)
\begin{itemize}[nosep]
  \item $\vec{P} = m\vec{g}$, sempre vertical para baixo,
        $g \approx 10\ \text{m/s}^2$.
  \item \textbf{Peso $\neq$ massa}: massa \'e em kg (quantidade de
        mat\'eria, a mesma em qualquer lugar); peso \'e em N (for\c{c}a,
        depende de $g$). Na Lua o peso diminui, a massa n\~ao.
  \item Conta r\'apida com a turma: pessoa de 60 kg $\Rightarrow$
        $P = 600$ N.
\end{itemize}

\medskip
\textbf{3. Normal} (5 min)
\begin{itemize}[nosep]
  \item For\c{c}a de contato da superf\'icie, sempre
        \textbf{perpendicular} \`a superf\'icie.
  \item Desenhar 3 casos lado a lado: livro na mesa (N para cima),
        quadro na parede (N horizontal), bloco na rampa
        (N perpendicular \`a rampa, \textbf{n\~ao} vertical!).
  \item Jogar a pergunta e deixar no ar: ``a normal \'e a rea\c{c}\~ao
        do peso?''. Responder s\'o na 3\textordfeminine{} Lei.
\end{itemize}

\medskip
\textbf{4. Atrito} (8 min)
\begin{itemize}[nosep]
  \item For\c{c}a de contato \textbf{paralela} \`a superf\'icie, contra
        o deslizamento (ou contra a tend\^encia de deslizar).
  \item F\'ormula: $f_{at} = \mu N$ (cin\'etico); no est\'atico vale
        $f_{at} \leq \mu_e N$, ou seja, o atrito est\'atico s\'o faz a
        for\c{c}a necess\'aria.
  \item Destacar que $\mu_e > \mu_c$: \'e mais dif\'icil come\c{c}ar a
        deslizar do que continuar deslizando. \'E assim que funciona o
        freio ABS, que j\'a caiu no ENEM (Quest\~ao 12 da lista).
  \item Lembrar que atrito n\~ao \'e sempre ruim: sem atrito n\~ao d\'a
        para andar. Isso liga com a 3\textordfeminine{} Lei e com a
        Quest\~ao 9, que vai ser feita em sala.
\end{itemize}

\medskip
\textbf{5. Tra\c{c}\~ao} (4 min)
\begin{itemize}[nosep]
  \item For\c{c}a transmitida por fio/corda/cabo esticado, sempre
        \textbf{puxando} na dire\c{c}\~ao do fio.
  \item Desenhar: lumin\'aria pendurada (T para cima, P para baixo,
        equil\'ibrio $\Rightarrow T = P$); corda de reboque.
\end{itemize}
\end{quadro}

\begin{perguntas}
\begin{itemize}[nosep]
  \item Depois do peso: ``a balan\c{c}a de farm\'acia mede massa ou
        peso?'' \textit{(Ela mede for\c{c}a --- peso/normal --- mas
        mostra em kg. Gera boa discuss\~ao.)}
  \item Depois do atrito: ``se n\~ao existisse atrito nenhum, o que
        acontece quando a pessoa tenta dar um passo?''
        \textit{(Desliza no lugar, n\~ao tem no que se apoiar.)}
\end{itemize}
\end{perguntas}

\begin{obs}
Usar uma cor de canet\~ao s\'o para os vetores de for\c{c}a em todos os
desenhos. Assim os alunos j\'a identificam as setas coloridas como
for\c{c}as, o que ajuda no Bloco 4.

N\~ao entrar em plano inclinado com decomposi\c{c}\~ao de vetores. S\'o
citar que a normal muda de dire\c{c}\~ao e seguir. Plano inclinado fica
para outra aula.
\end{obs}

\begin{corte}
Cortar a tra\c{c}\~ao (item 5), ela aparece de novo quando algum
exerc\'icio pedir. N\~ao cortar o atrito: tem 3 quest\~oes de atrito na
lista e \'e o assunto de din\^amica que mais cai no ENEM.
\end{corte}

% ===================================================================
\bloco{3}{3\textordfeminine{} Lei --- A\c{c}\~ao e Rea\c{c}\~ao}{15}
% ===================================================================

\objetivo{enunciar a 3\textordfeminine{} Lei e deixar claro que o par
a\c{c}\~ao--rea\c{c}\~ao atua em corpos diferentes, que \'e onde os
alunos mais erram. Resolver tamb\'em a d\'uvida da normal que ficou do
Bloco 2.}

\begin{quadro}
\textbf{1. Enunciado} (4 min)
\begin{itemize}[nosep]
  \item Escrever na faixa de resumo:
  \begin{center}
    \textbf{3\textordfeminine{} Lei:} se A empurra B com $\vec{F}$,
    B empurra A com $-\vec{F}$.\\
    \small (mesma intensidade, mesma dire\c{c}\~ao, sentidos opostos,
    \textbf{corpos diferentes})
  \end{center}
  \item Sublinhar \textbf{corpos diferentes} duas vezes. \'E por isso
        que a\c{c}\~ao e rea\c{c}\~ao \textbf{nunca se cancelam}: cada
        uma age em um corpo.
\end{itemize}

\medskip
\textbf{2. Tr\^es exemplos no quadro} (6 min)
\begin{itemize}[nosep]
  \item \textbf{Empurrar a parede} (exemplo do planejamento original):
        m\~ao empurra parede $\rightarrow$, parede empurra m\~ao
        $\leftarrow$ (de patins, a pessoa recua).
  \item \textbf{Caminhar}: p\'e empurra o ch\~ao para tr\'as, ch\~ao
        (atrito!) empurra a pessoa para a frente. Quem anda \'e movido
        pelo ch\~ao. (Quest\~ao 9 da lista, ENEM 2013, vai ser feita em
        sala.)
  \item \textbf{Foguete}: expele g\'as para baixo, g\'as empurra o
        foguete para cima. Funciona at\'e no v\'acuo, n\~ao precisa de
        ``apoio no ar'' (Quest\~ao 11 da lista, UTFPR 2023).
\end{itemize}

\medskip
\textbf{3. Resolver a d\'uvida da normal} (5 min)
\begin{itemize}[nosep]
  \item Livro na mesa: desenhar \textbf{dois} diagramas separados
        (livro / mesa).
  \item Peso do livro = Terra puxa livro. Rea\c{c}\~ao: \textbf{livro
        puxa a Terra} (age na Terra!).
  \item Normal = mesa empurra livro. Rea\c{c}\~ao: livro empurra a
        mesa.
  \item Concluir no quadro: \textbf{peso e normal N\~AO s\~ao par
        a\c{c}\~ao--rea\c{c}\~ao} (agem no mesmo corpo e t\^em
        naturezas diferentes). Eles s\'o se equilibram nesse caso.
\end{itemize}
\end{quadro}

\begin{perguntas}
\begin{itemize}[nosep]
  \item ``Se a for\c{c}a que o cavalo faz na carro\c{c}a \'e igual \`a
        que a carro\c{c}a faz no cavalo, como a carro\c{c}a sai do
        lugar?''
        \textit{(Cada for\c{c}a age em um corpo. A carro\c{c}a acelera
        porque a resultante NELA --- tra\c{c}\~ao menos atrito --- \'e
        positiva.)}
  \item ``A Terra puxa uma pessoa com 600 N. Com quanto a pessoa puxa a
        Terra?''
        \textit{(600 N tamb\'em. S\'o que a acelera\c{c}\~ao da Terra
        \'e desprez\'ivel, porque a massa dela \'e enorme ---
        2\textordfeminine{} Lei.)}
\end{itemize}
\end{perguntas}

\begin{obs}
A armadilha cl\'assica do ENEM \'e ``a\c{c}\~ao e rea\c{c}\~ao se
anulam, logo nada se move''. \'E exatamente o argumento do garoto na
Quest\~ao 10 da lista (ENEM PPL 2012), que vai ser feita em sala. Se a
discuss\~ao do cavalo j\'a tiver aparecido aqui, essa quest\~ao fica
f\'acil.
\end{obs}

\begin{corte}
Cortar o exemplo do foguete (a Quest\~ao 11 cobre a mesma ideia como
tarefa). N\~ao cortar a parte da normal, que \'e a mais importante do
bloco.
\end{corte}

% ===================================================================
\bloco{4}{Aplica\c{c}\~ao: Diagrama de For\c{c}as + Exemplo}{15}
% ===================================================================

\objetivo{transformar as tr\^es leis em m\'etodo (isolar o corpo,
desenhar as for\c{c}as, aplicar $F_R = ma$) e resolver um exemplo
num\'erico completo no quadro.}

\begin{quadro}
\textbf{1. Passo a passo (escrever no quadro e numerar)} (4 min)
\begin{enumerate}[nosep]
  \item Isolar o corpo (desenhar s\'o ele);
  \item Desenhar \textbf{todas} as for\c{c}as que agem \textbf{nele}
        (peso? contato? fio?);
  \item Escolher o eixo positivo no sentido do movimento;
  \item Aplicar $F_R = ma$ em cada eixo.
\end{enumerate}

\medskip
\textbf{2. Treino r\'apido de diagrama (a turma vai ditando)} (4 min)
\begin{itemize}[nosep]
  \item Livro parado na mesa: $\vec{P}$, $\vec{N}$.
  \item Bloco puxado por corda no ch\~ao com atrito:
        $\vec{P}$, $\vec{N}$, $\vec{T}$, $\vec{f}_{at}$.
  \item Pergunta-armadilha: ``e a `for\c{c}a do movimento'?''
        \textit{(N\~ao existe. Movimento n\~ao \'e for\c{c}a.)}
\end{itemize}

\medskip
\textbf{3. Exemplo resolvido (estilo UTFPR)} (7 min)
\begin{itemize}[nosep]
  \item Ditar o enunciado: bloco de $4{,}0$ kg no ch\~ao, puxado na
        horizontal com $F = 16{,}0$ N, atrito desprez\'ivel.
        (a) acelera\c{c}\~ao? (b) velocidade depois de 5 s, partindo de
        $18$ km/h?
  \item Resolver no quadro seguindo o passo a passo:
  \[
    a = \frac{F_R}{m} = \frac{16}{4} = 4\ \text{m/s}^2
  \]
  \[
    18\ \text{km/h} = 5\ \text{m/s} \qquad
    v = v_0 + at = 5 + 4\cdot 5 = 25\ \text{m/s}
  \]
  \item Comentar que esse exemplo \'e a Quest\~ao 4 da lista, que caiu
        exatamente assim na UTFPR Ver\~ao 2025, e que ela vai ser
        refeita completa no Bloco 5 (com o item do espa\c{c}o
        percorrido).
\end{itemize}
\end{quadro}

\begin{obs}
Revisar r\'apido a convers\~ao km/h $\to$ m/s (dividir por 3,6) aqui,
para n\~ao travar o Bloco 5.
\end{obs}

\begin{corte}
Se o Bloco 2 estourou o tempo: pular o item 2 (treino de diagramas) e
ir direto para o exemplo resolvido. Os diagramas aparecem de novo nas
quest\~oes de sala.
\end{corte}

% ===================================================================
\bloco{5}{Quest\~oes ENEM/UTFPR ao vivo}{25}
% ===================================================================

\objetivo{terminar a aula resolvendo quest\~oes de verdade do ENEM e da
UTFPR, projetadas uma por uma e resolvidas ao vivo no quadro. \'Unico
bloco que usa projetor. As quest\~oes est\~ao no arquivo
\texttt{lista\_questoes\_dinamica.pdf} (as marcadas com
{\scshape\color{vinho}[fazer em sala]}).}

\begin{projetor}
Ligar o projetor (primeira vez na aula) com o arquivo
\texttt{lista\_questoes\_dinamica.pdf}.

\smallskip
\textbf{Para cada quest\~ao:}
\begin{enumerate}[nosep]
  \item Projetar o enunciado e ler em voz alta. Dar uns 2 min para
        tentarem sozinhos.
  \item Votar: ``quem marcou A? B?...'' (anotar no canto do quadro).
  \item Resolver \textbf{no quadro}, come\c{c}ando pelo diagrama de
        for\c{c}as quando fizer sentido.
  \item Fechar falando \textbf{qual lei} a quest\~ao cobrou (e marcar
        na faixa de resumo).
\end{enumerate}
\end{projetor}

\begin{quadro}
\textbf{Roteiro de quest\~oes} (n\'umeros da lista):

\medskip
\begin{tabular}{@{}clll@{}}
  \toprule
  \# & Fonte & Cobra & Tempo \\
  \midrule
  Q1  & ENEM PPL 2012 & 1\textordfeminine{} Lei (bola no vag\~ao)   & 4 min \\
  Q10 & ENEM PPL 2012 & 3\textordfeminine{} Lei (m\~ae e o m\'ovel) & 3 min \\
  Q9  & ENEM 2013     & 3\textordfeminine{} Lei + atrito (rampa)    & 4 min \\
  Q4  & UTFPR Ver\~ao 2025 & 2\textordfeminine{} Lei num\'erica     & 6 min \\
  Q5  & ENEM 2017     & 2\textordfeminine{} Lei + gr\'afico (cinto) & 5 min \\
  \midrule
  Q6  & UTFPR Ver\~ao 2026 & Frenagem (b\^onus, se der tempo)       & 5 min \\
  \bottomrule
\end{tabular}

\medskip
A ordem \'e essa de prop\'osito: come\c{c}ar pelas duas conceituais,
que s\~ao r\'apidas, depois a cl\'assica do atrito ao andar, depois a
num\'erica da UTFPR (que \'e o exemplo do Bloco 4) e por \'ultimo a do
gr\'afico, que \'e o estilo t\'ipico do ENEM.
\end{quadro}

\begin{obs}
Melhor fazer 5 quest\~oes com calma do que 6 correndo. No ENEM o erro
mais comum \'e de leitura, n\~ao de conta, ent\~ao vale gastar tempo no
enunciado.
\end{obs}

\begin{corte}
Menos de 15 min? Fazer Q1, Q9 e Q4, que cobrem as tr\^es leis.
Menos de 10? Q9 e Q4.
\end{corte}

% ===================================================================
\bloco{6}{Encerramento e Tarefa}{10}
% ===================================================================

\begin{quadro}
\textbf{1. Resumo} (5 min) --- conferir a faixa de resumo do quadro e
completar a coluna de exemplos com a turma ditando:

\medskip
\begin{tabular}{@{}lll@{}}
  \toprule
  Lei & Frase-chave & Exemplo \\
  \midrule
  1\textordfeminine{} & $F_R=0 \Rightarrow$ repouso ou MRU & passageiro sem cinto \\
  2\textordfeminine{} & $F_R = ma$                          & carrinho vazio/cheio \\
  3\textordfeminine{} & $\vec{F}_{AB} = -\vec{F}_{BA}$, corpos diferentes & andar, foguete \\
  \bottomrule
\end{tabular}

\medskip
\textbf{2. Tarefa} (3 min) --- distribuir/enviar a lista de quest\~oes:
\begin{itemize}[nosep]
  \item Resolver as quest\~oes que sobraram da lista (as sem marca de
        sala);
  \item Come\c{c}ar por Q2, Q3 e Q11 (uma de cada lei);
  \item Q15 e Q16 s\~ao desafios. Vale tentar; elas voltam nas
        pr\'oximas aulas (energia e colis\~oes).
\end{itemize}
\end{quadro}

\begin{perguntas}
Para fechar:
\begin{itemize}[nosep]
  \item ``Qual das 3 leis voc\^es acham que mais cai?''
        \textit{(Todas caem. A 1\textordfeminine{} e a
        3\textordfeminine{} caem como conceito/pegadinha e a
        2\textordfeminine{} como conta.)}
  \item ``Qual pergunta do come\c{c}o da aula a gente respondeu hoje?''
\end{itemize}
\end{perguntas}

\begin{obs}
Avisar qual \'e a pr\'oxima aula (conforme o plano dos 5 encontros) e
falar que plano inclinado e sistemas de blocos v\~ao aparecer l\'a,
para ningu\'em achar que esse conte\'udo ficou faltando.
\end{obs}

% ===================================================================
\secao{Adapta\c{c}\~ao por Tempo Dispon\'ivel}
% ===================================================================

{\renewcommand{\arraystretch}{1.5}
\begin{tabularx}{\linewidth}{@{}lX@{}}
  \toprule
  \textbf{Tempo} & \textbf{O que fazer} \\
  \midrule
  $\approx 1$h30 &
    B0 (5 min) + B1 (12 min) + B2 sem tra\c{c}\~ao (25 min) +
    B3 (12 min) + B4 s\'o exemplo (8 min) + 4 quest\~oes (20 min) +
    fechamento (5 min) \\
  $\approx 2$h &
    Plano completo como descrito \\
  $\approx 2$h30 &
    Plano completo + Q6 garantida + mais treino de diagramas no B4 +
    discutir a Q11 (astronauta) em sala \\
  \bottomrule
\end{tabularx}}

\bigskip

{\scshape\color{vinho}prioridade dos t\'opicos para enem/utfpr}
\begin{enumerate}[nosep]
  \item \textbf{3\textordfeminine{} Lei como conceito} (pares em corpos
        diferentes, atrito ao andar) --- o ENEM cobra muito isso
  \item \textbf{Atrito} (est\'atico vs.\ cin\'etico, ABS)
  \item \textbf{2\textordfeminine{} Lei num\'erica} ($F=ma$ +
        cinem\'atica) --- estilo UTFPR
  \item \textbf{1\textordfeminine{} Lei / in\'ercia} (frenagens, curvas)
  \item \textbf{Peso $\times$ massa} e a d\'uvida peso--normal
\end{enumerate}

\end{document}
